Nog in te vullen
\section{Interview Hans Kok - ABN AMRO}
Senior z/OS Specialist - interview fysiek en online
\subsubsection{interview 1}
Vraag: Welke concepten of vaardigheden vond u het moeilijkst om te leren
bij de start van uw carrière?

Antwoord: In het begin weet je niks, en dat gold voor ons allemaal. We hadden VM
waar we de hele dag met elkaar chatten (we werden verdeeld over drie
locaties, de VM-en waren verbonden), het was erg moeilijk om iets te
betekenen.
Eerste opdracht die ik kreeg was het debuggen van een PL/I
programmaatje, eigenlijk waren het 3 programmaatjes, het eerste
programmaatje bouwde een table op met tijden en system commando’s,
programma twee pakte dat op en gaf dat door aan programma 3 (assembler)
die voerde de systeemcommando’s uit.
Het programma kende wat vreemde abends en kwa commando’s liep hij uit
zijn tijd, het kostte me enkele weken, maar het was leerzaam. De abends
kwamen door een te klein gedefinieerde table, het uit de tijd lopen kwam
door vertraging bij het inswappen, er was geen terugkoppelingcheck
ingebouwd.
Ik had een mentor waar ik niet zo veel aan had. Naast me zat een
Amerikaan op assignment die altijd en met veel geduld antwoord gaf op
mijn vragen, daar heb ik veel aan gehad.
Toen ik verantwoordelijk werd voor het afhandelen van incidenten kreeg
ik het idee dat ik begon te functioneren, het heeft zo’n twee jaar
geduurd dat ik dat gevoel kreeg. Alles wat misging op een stuk of 8
lpars kwam bij mij terecht, en je moest de systemen en de boeken
induiken om het op te lossen, daar heb ik het meest van geleerd.
Dat zou dus een tip zijn aan beginners, zorg dat je bij de club komt die
incidenten afhandelt, en gebruik de middelen die tot je beschikking
staan, syslog, netlog en msgid’s.

Vraag: Welke kennis of competenties waren essentieel bij het werken binnen
uw organisatie?

Antwoord: Zelf heb ik in 13 bedrijven gewerkt. Als je binnenkomt in een bedrijf
ken je de mensen niet. Naast een degelijke kennis van zaken is het
belangrijk om goede contacten te hebben binnen verschillende afdelingen.
Iets geregeld krijgen werkt beter als je de mensen kent, zo moet je zelf
ook open staan, werkt gewoon sneller en efficienter dan wanneer je via
officiele kanalen mensen benadert. Daar moet natuurlijk bij  opgemerkt
worden dat je je wel altijd moet conformeren aan het change process, het
gaat om kleine onschuldige wijzigingen.

Vraag: Welk advies zou u geven aan een student die een carrière in
mainframe sector?

Antwoord: Met adviezen moet je altijd voorzichtig zijn, wat vandaag beroepsmatig
veel gevraagd wordt kan over enkele jaren overbodig zijn en andersom.
Op mijn gevoel zou ik zeggen stort je in een carriere in de mainframe
sector, maar blijf vooral goed op de hoogte van wat er allemaal speelt.
Ik verbaas me over de mogelijkheden van AI en kan niet inschatten in
hoeverre dat een reductie zal zijn van de vraag binnen ons vakgebied.
Door de jaren heen heb ik ervaren dat ons vakgebied onderschat wordt, we
zijn maar knoppendrukkers is een idee wat heerst bij een deel van het
hogere echlon, nou nee, ons werk vraagt een academisch denkniveau.

Vraag: Denk je dat de website die ik ontwikkel een nuttig hulpmiddel zou
kunnen zijn?

Antwoord: Dat denk ik absoluut. Wat je nu geproduceerd hebt vind ik op een aantal
punten al verhelderend als ik terug denk aan de verwarring in de eerste
periode van onze opleiding.
Het is vooral belangrijk dat jonge mensen enthousiast gemaakt worden
voor dit vakgebied, voor mij bestaat er geen mooier vak binnen de ICT
dan z/OS Syspro of Sysadmin.
\subsubsection{interview 2}

\section{Interview Gorkem Cekmen - IsBank}
senior z/OS sys prog bij isbank - interview fysiek op MTE (interview afgenomen in het Turks en vertaald naar Nederlands)

Vraag: Wat vond u het moeilijkst bij de start van uw mainframecarrière?

Antwoord Gorkem: In het begin was het vooral moeilijk om z/OS als geheel te begrijpen. Het is geen eenvoudige technologie waar je snel doorheen bent. Alles hangt met elkaar samen: het operating system, de subsystems, storage, security… Je moet echt begrijpen hoe die componenten samenwerken. In het begin zijn dat vrij abstracte concepten en dat maakt het moeilijk. Je ziet wel losse onderdelen, maar het grotere geheel begrijpen vraagt tijd.

Vraag: Hoe hebt u die moeilijkheden aangepakt?

Antwoord Gorkem: Ik had het geluk dat ik een collega had met meer dan 20 jaar ervaring. Hij heeft mij eigenlijk stap voor stap opgeleid. Veel kennis in de mainframewereld wordt nog altijd doorgegeven via ervaring en mentoring. Daarnaast zocht ik ook zelf naar informatie,  op mainframe forums. En als ik vastzat, belde ik gewoon collega’s. Het is een combinatie van zelfstudie en hulp vragen aan ervaren mensen.

{Vraag: Welke technologie of vaardigheid binnen z/OS vindt u het moeilijkst?} \\

{Antwoord Gorkem:} Voor mij is dat zonder twijfel assembler. Je moet een goed begrip hebben van hoe het systeem intern functioneert, zoals registers, geheugenbeheer en instructiesets. In mijn team (Isbank) moet je ook assembler kunnen en in het begin had ik veel moelijk mee. Kleine fouten kunnen grote gevolgen hebben, en debugging is niet altijd eenvoudig. Het is iets dat je enkel echt onder de knie krijgt door veel ervaring en praktijk.

Vraag: Welk advies zou u geven aan studenten die een carrière in de mainframe-sector overwegen?

Antwoord Gorkem: Probeer eerst de basis goed te begrijpen en wees niet bang om vragen te stellen. Het is een complexe omgeving, dus het is normaal dat je niet alles meteen begrijpt. Zoek actief hulp, leer van ervaren collega’s en blijf oefenen. Met tijd en ervaring wordt alles duidelijker.


\section{Interview Wim Van Weven - ABN AMRO}
modernisatie (liberty, websphere) - interview fysiek



\section{Interview Jolanda Berkhout - ABN AMRO}
product owner/cobol dev - interview fysiek



\section{Interview Brahm Landrechts - ABN AMRO}
IMS specialist bij abn amro - interview fysiek



\section{Interview Herman Hartman - Planet NL}
z/OS sysadmin/sysprog - interview online

Vraag: Wat is uw algemene indruk van het voorgestelde competentieframework? (Versie van de bachelorproef en de website in maart)

Antwoord Herman: De inventory geeft de indruk van een eerste “brain dump”. Dat is op zich niet slecht als startpunt, maar het roept wel de vraag op waar de topics precies vandaan komen. Je maakt al een onderscheid tussen System Administrator level 1 en level 2, maar dat lijkt daar een beetje los van te staan. Het zou sterker zijn als de structuur eerst duidelijk wordt bepaald en daarna pas ingevuld.

Vraag: Hoe zou u het framework beter structureren?

Antwoord Herman: Ik zou vertrekken van meer neutrale functionele domeinen. Bijvoorbeeld: hardware, hypervisor, operating system, system software, application, operating, security, en governance \& compliance. Binnen elk van die domeinen kan een system administrator dan zijn taken uitvoeren of services leveren. Als je eerst die structuur hebt, wordt het ook makkelijker om te mappen naar bestaande frameworks.

Vraag: Hoe ziet u de rol van een System Administrator in de praktijk?

Antwoord Herman: In een gemiddelde z/OS-omgeving heb je een vrij grote IT-organisatie. Er is veel specialisatie: operations, security, incident management, problem management, user support, capacity, reporting, compliance. Dat sluit ook aan bij frameworks zoals ITIL, ASL en BiSL. Het is dus belangrijk dat je framework die realiteit weerspiegelt en niet alles als één rol bekijkt.

Vraag: Zijn er volgens u belangrijke onderdelen die nog ontbreken?

Antwoord Herman: Ja, ik mis nog een stuk rond system development en de interactie met system administrators. Daarnaast zou knowledge transfer ook een plaats mogen krijgen, bijvoorbeeld via een soort sempai/kohai-model, waarbij ervaren mensen hun kennis doorgeven aan juniors.

Vraag: Wat vindt u van de huidige opleidingen en frameworks?

Antwoord Herman: Het kan nuttig zijn om expliciet te maken waar bestaande trainingen en frameworks tekortschieten. IBM heeft daar al een goed overzicht van gegeven in recente publicaties. Het is interessant om daarop verder te bouwen, maar ook aan te tonen wat er nog ontbreekt.

Vraag: Heeft u concrete feedback op de Excel-matrix? (Een excel sheet met alle onderwerpen dat ik gebundeld heb)

Antwoord Herman: Ja, enkele praktische dingen. Bij hardware zou je bijvoorbeeld een startup-component kunnen toevoegen. Bij JCL zag ik een kleine fout: IEFB14 moet IEFBR14 zijn, en utilities zou je beter als aparte lijn zetten. Security zou ik eerder “Authorisation” noemen, met RACF als implementatie. Verder kun je zaken toevoegen zoals SNA (historisch), Db2 en IMS apart vermelden, en batch automation tools zoals Tivoli Workload Scheduler. Ook een aparte lijn “Operating z/OS” met SMF, RMF en UNIX shell lijkt me nuttig.

Vraag: Welke concepten of vaardigheden vond u het moeilijkst om te leren bij de start van uw carrière?

Antwoord Herman: De grootste uitdaging is het begrijpen van het geheel. z/OS is geen losse technologie, maar een ecosysteem van componenten die sterk met elkaar verbonden zijn. In het begin is het moeilijk om te zien hoe alles samenwerkt: operating system, subsystems, security, networking. Daarnaast zijn zaken zoals JCL en debugging in zo’n omgeving niet intuïtief als je van andere platformen komt. Dat leer je vooral door ervaring en door samen te werken met collega’s.

Vraag: Welk advies zou u geven aan een student die een carrière in de mainframe-sector overweegt?

Antwoord Herman: Focus eerst op de fundamenten en probeer te begrijpen hoe het systeem in elkaar zit, in plaats van enkel tools te leren. Daarnaast is hands-on ervaring cruciaal. Je leert dit soort technologie niet alleen uit boeken. Werk samen met ervaren collega’s en stel vragen. Laat je ook niet afschrikken door de complexiteit die is er, maar net daardoor is de vraag naar goede profielen groot. Als je erin investeert, kun je een zeer stabiele en interessante carrière opbouwen.

\section{Interview Rudie Claemans - ABN AMRO}
senior z/OS - interview fysiek


\section{Interview Peter Heijlo - ABN AMRO}
mainframe infrastructure team manager - interview fysiek 
\section{Interview Murat Kilickaya - Broadcom}
broadcom representative. over softskills bij junior mainframe - interview fysiek op MTE
\section{Interview Berkay Ariturk - Halkbank}
junior z/OS sysprog bij halkbank - interview online 
\section{Interview Jean Rollet - (Hercules gebruiker)}

Jean Dupont is actief binnen de Hercules-community en gebruikt z/OS-emulatie in een leer- en experimenteercontext. Het gesprek vond online plaats en focuste op het gebruik van emulators binnen de competentieontwikkeling van beginnende system programmers en system administrators.

{Vraag: Beschikt u over een persoonlijke z/OS-omgeving om autonoom te werken?} \\
{Antwoord Jean:} Ja, ik werk met een eigen omgeving gebaseerd op de Hercules-emulator. Dat laat mij toe om z/OS lokaal te draaien op een gewone machine. Het is uiteraard geen productieomgeving zoals in bedrijven, maar voor leerdoeleinden en experimenten is het zeer waardevol. Je kan er dingen testen zonder risico, wat in een echte mainframe-omgeving vaak niet mogelijk is.

{Vraag: Hoe belangrijk zijn emulators volgens u voor beginners in de mainframewereld?} \\
{Antwoord Jean:} Ik denk dat ze echt handig kunnen zijn. De grootste drempel in de mainframewereld is toegang. Niet iedereen heeft meteen toegang tot een echte z/OS-omgeving binnen een bedrijf. Met een emulator zoals Hercules kan je die drempel verlagen. Studenten kunnen zelfstandig oefenen, fouten maken en opnieuw proberen. Dat versnelt het leerproces aanzienlijk.

{Vraag: Welke vaardigheden kunnen studenten ontwikkelen met zo’n omgeving?} \\
{Antwoord Jean:} Je hebt meer vreede dan zXplore platform van IBM. Je kan experimenteren met IPL, onderliggende subsystemen en dergelijke.

{Vraag: Kan een emulator ook nuttig zijn voor meer gevorderde competenties?} \\
{Antwoord Jean:} Zeker. Ik gebruik Hercules in combinatie met machine learning toepassingen. Enerzijds om bestaande systemen te analyseren, zoals legacy code en systeemstructuren. Anderzijds om scenario’s te testen rond migratie naar open source oplossingen. Het is een veilige sandbox waarin je complexe ideeën kan uitwerken zonder impact op productie.

{Vraag: Wat vindt u van de onderwerpen in de Excel-sheet die ik heb opgesteld voor het z/OS competentieframework?} \\
{Antwoord Jean:} Ik vind het een sterke basis, vooral omdat je duidelijk probeert om het volledige landschap van z/OS te capteren. Je raakt veel relevante technologieën en concepten aan maar je moet eigenlijk ook deftig kunnen voorstellen aan de gebruiker van je competentieframework. 
Ik zou ook IPL procedure toevoegen stap per stap, dan kan jij ook een praktisch oefening toevoegen van IPL via emulator. Maar in het algemeen: het is een goed vertrekpunt dat vooral nog wat verfijning nodig heeft richting praktische toepasbaarheid.